\documentclass[man,12pt]{apa7}

\usepackage[spanish,es-tabla]{babel}
\usepackage[utf8]{inputenc}
\usepackage[T1]{fontenc}
\usepackage{csquotes}
\usepackage{graphicx}
\usepackage{booktabs}
\usepackage{float}
\usepackage{xcolor}
\usepackage{listings}
\usepackage{hyperref}
\usepackage[
  style=apa,
  sortcites=true,
  sorting=nyt,
  backend=biber
]{biblatex}

\addbibresource{referencias.bib}

\hypersetup{
  colorlinks=true,
  linkcolor=black,
  urlcolor=blue,
  citecolor=black
}

\lstdefinestyle{codigo}{
  basicstyle=\ttfamily\small,
  breaklines=true,
  frame=single,
  backgroundcolor=\color{gray!8},
  rulecolor=\color{gray!35}
}

\title{Evaluacion de Calidad de Software del Proyecto Mentes Creativas: Ciencias Naturales}
\shorttitle{Calidad de Software en Mentes Creativas}
\author{Nombre del integrante 1 \\ Nombre del integrante 2 \\ Nombre del integrante 3}
\affiliation{Universidad Cooperativa de Colombia}
\course{Calidad de Software}
\professor{Mg. Gustavo Sanchez Rodriguez}
\duedate{Fecha de entrega}

\begin{document}
\maketitle

\begin{abstract}
Este documento presenta la evaluacion de calidad del proyecto \textit{Mentes Creativas: Ciencias Naturales}, una aplicacion frontend construida con React y Vite para apoyar procesos de enseñanza-aprendizaje en estudiantes de cuarto y quinto grado. El proyecto desarrolla tres temas: ecosistemas, ciclo del agua y cuerpo humano. La evaluacion se fundamenta en la norma ISO/IEC 25010, con enfasis en usabilidad, capacidad de aprendizaje y operabilidad. Tambien se documentan pruebas unitarias con Jest, integracion con Postman, pruebas de sistema con JMeter, integracion continua con GitHub Actions y despliegue del frontend en Vercel.
\end{abstract}

\keywords{calidad de software, ISO/IEC 25010, React, Jest, Vercel, ciencias naturales}

\section{Introduccion}

La calidad del software permite verificar que una aplicacion cumpla los requisitos funcionales, tecnicos y de experiencia definidos para sus usuarios. En este trabajo se desarrolla y evalua una aplicacion educativa multimedia para el Colegio Mentes Creativas, orientada a estudiantes de cuarto y quinto grado en el area de Ciencias Naturales.

El proyecto se implementa como un frontend en React con Vite. Su proposito es facilitar el aprendizaje ludico mediante recursos visuales, audio por sintesis de voz, animaciones y modelos graficos 3D basados en CSS. La evaluacion se realiza a partir del modelo ISO/IEC 25010, pruebas automatizadas, pruebas de integracion, pruebas de carga, validacion de despliegue y criterios de aceptacion \parencite{iso25010,pressman2020}.

\section{Escenario Asignado y Justificacion}

El Colegio Mentes Creativas busca implementar una aplicacion multimedia que apoye procesos de enseñanza-aprendizaje en areas clave del curriculo. Para esta entrega se selecciono el area de Ciencias Naturales, debido a que permite explicar fenomenos observables mediante recursos interactivos y representaciones visuales.

La aplicacion desarrolla tres temas:

\begin{itemize}
  \item Ecosistemas.
  \item Ciclo del agua.
  \item Cuerpo humano.
\end{itemize}

La seleccion de estos temas permite trabajar conceptos fundamentales del grado cuarto y quinto, como relaciones entre seres vivos, cambios del agua en la naturaleza y sistemas del cuerpo humano. El uso de componentes interactivos busca favorecer la comprension, la exploracion autonoma y la retroalimentacion inmediata.

\section{Requerimientos Funcionales}

La Tabla \ref{tab:requerimientos} presenta los requerimientos funcionales definidos para el proyecto.

\begin{table}[H]
\centering
\caption{Requerimientos funcionales del proyecto}
\label{tab:requerimientos}
\begin{tabular}{p{2.2cm}p{10.2cm}}
\toprule
\textbf{ID} & \textbf{Requerimiento funcional} \\
\midrule
RF-01 & El sistema debe permitir seleccionar entre los temas Ecosistemas, Ciclo del agua y Cuerpo humano. \\
RF-02 & El sistema debe mostrar objetivos de aprendizaje por cada tema. \\
RF-03 & El sistema debe presentar palabras clave y conceptos principales de cada tema. \\
RF-04 & El sistema debe incluir un recurso multimedia interactivo por tema. \\
RF-05 & El sistema debe permitir reproducir una guia auditiva mediante sintesis de voz. \\
RF-06 & El sistema debe incluir una actividad de laboratorio guiada por cada tema. \\
RF-07 & El sistema debe presentar una evaluacion corta con retroalimentacion inmediata. \\
RF-08 & El sistema debe publicar un endpoint estatico con los temas disponibles para pruebas de integracion. \\
\bottomrule
\end{tabular}
\end{table}

\section{Norma y Modelo de Calidad Aplicado}

La norma ISO/IEC 25010 define un modelo de calidad del producto software que agrupa caracteristicas como adecuacion funcional, eficiencia de desempeño, compatibilidad, usabilidad, fiabilidad, seguridad, mantenibilidad y portabilidad \parencite{iso25010}.

Para este proyecto se selecciona la caracteristica \textbf{usabilidad}, porque el usuario principal es un estudiante de cuarto o quinto grado. En este contexto, la calidad no depende solo de que la aplicacion funcione, sino de que sea comprensible, operable y adecuada para apoyar el aprendizaje.

Los subatributos seleccionados son:

\begin{itemize}
  \item \textbf{Capacidad de aprendizaje}: facilidad con la que el estudiante comprende como usar el sistema y aprende el contenido.
  \item \textbf{Operabilidad}: facilidad para navegar, seleccionar temas, activar recursos y responder evaluaciones.
\end{itemize}

\section{Metricas de Calidad Definidas}

\begin{table}[H]
\centering
\caption{Metricas asociadas a ISO/IEC 25010}
\label{tab:metricas}
\begin{tabular}{p{3cm}p{4.2cm}p{4.8cm}}
\toprule
\textbf{Subatributo} & \textbf{Metrica} & \textbf{Criterio de aceptacion} \\
\midrule
Capacidad de aprendizaje & Finalizacion de una evaluacion corta por tema. & El estudiante responde al menos una pregunta y recibe retroalimentacion visible. \\
Capacidad de aprendizaje & Claridad de objetivos y palabras clave. & Cada tema contiene minimo tres objetivos y cinco palabras clave. \\
Operabilidad & Navegacion entre temas y secciones. & El usuario puede cambiar de tema y abrir Exploracion, Laboratorio y Evaluacion sin errores. \\
Operabilidad & Rendimiento percibido. & El sitio desplegado responde con promedio menor a cinco segundos en JMeter. \\
\bottomrule
\end{tabular}
\end{table}

\section{Arquitectura Tecnica}

El proyecto se implementa con React, Vite y TypeScript. La interfaz se divide en componentes reutilizables: un menu lateral de temas, un lienzo visual para recursos 3D, una seccion de exploracion, un laboratorio guiado y un panel de evaluacion.

La estructura principal del repositorio es:

\begin{lstlisting}[style=codigo]
src/
  components/
  data/
  test/
public/
  api/
docs/
  evidencias/
  jmeter/
  overleaf/
  postman/
.github/workflows/
\end{lstlisting}

\section{Proceso de Despliegue en Vercel}

El despliegue se realiza conectando el repositorio de GitHub con Vercel. La configuracion usada es:

\begin{itemize}
  \item Framework preset: Vite.
  \item Install command: \texttt{npm install}.
  \item Build command: \texttt{npm run build}.
  \item Output directory: \texttt{dist}.
\end{itemize}

La URL publica del despliegue es:

\begin{quote}
\textbf{URL de Vercel:} \url{https://mentes-creativas-ciencias-naturales-rho.vercel.app}
\end{quote}

\textbf{Evidencia requerida:} agregar captura del despliegue exitoso y de la aplicacion funcionando en produccion.

\section{Pruebas Unitarias Implementadas}

Las pruebas unitarias se implementan con Jest y React Testing Library \parencite{jestDocs,testingLibraryDocs}. Estas pruebas validan:

\begin{itemize}
  \item Renderizado del titulo y los tres temas.
  \item Cambio de tema desde el menu.
  \item Resolucion de preguntas con retroalimentacion.
  \item Activacion del recurso de audio mediante Web Speech API.
  \item Consistencia de los datos educativos.
\end{itemize}

El comando de ejecucion es:

\begin{lstlisting}[style=codigo]
npm test
\end{lstlisting}

\textbf{Evidencia requerida:} agregar captura de la terminal con pruebas exitosas.

\section{Automatizacion de Integracion Continua}

Se configura GitHub Actions mediante el archivo \texttt{.github/workflows/ci.yml}. El flujo ejecuta instalacion de dependencias, pruebas unitarias y construccion del proyecto en cada push o pull request hacia la rama \texttt{main}.

\begin{lstlisting}[style=codigo]
npm install
npm test
npm run build
\end{lstlisting}

La integracion continua permite detectar errores antes del despliegue y aporta evidencia objetiva del proceso de aseguramiento de calidad \parencite{githubActionsDocs}.

\section{Pruebas de Integracion con Postman}

La coleccion de Postman se encuentra en:

\begin{lstlisting}[style=codigo]
docs/postman/mentes-creativas.postman_collection.json
\end{lstlisting}

Los casos incluidos son:

\begin{itemize}
  \item Validar carga de la pagina principal.
  \item Validar estado 200 del endpoint \texttt{/api/lessons.json}.
  \item Verificar que el endpoint contiene tres temas.
  \item Confirmar que uno de los temas es Ciclo del agua.
\end{itemize}

\textbf{Evidencia requerida:} agregar capturas de Postman con pruebas exitosas.

\section{Pruebas de Sistema con JMeter}

La prueba de carga se diseña en JMeter para simular 25 usuarios concurrentes, con dos iteraciones por usuario, sobre la URL publica de Vercel. El archivo se encuentra en:

\begin{lstlisting}[style=codigo]
docs/jmeter/mentes-creativas-carga.jmx
\end{lstlisting}

Los endpoints evaluados son:

\begin{itemize}
  \item \texttt{/}
  \item \texttt{/api/lessons.json}
\end{itemize}

El criterio de aceptacion definido es mantener un tiempo promedio menor a cinco segundos y una tasa de error cercana a cero.

\section{Pruebas de Implantacion}

Las pruebas de implantacion validan que el frontend se encuentre disponible en Vercel y que las funcionalidades principales se ejecuten correctamente en produccion. Se debe verificar:

\begin{itemize}
  \item Carga inicial del sitio.
  \item Navegacion entre los tres temas.
  \item Acceso a exploracion, laboratorio y evaluacion.
  \item Endpoint \texttt{/api/lessons.json}.
  \item Consola del navegador sin errores visibles.
\end{itemize}

\section{Checklist de Aceptacion}

\begin{table}[H]
\centering
\caption{Checklist de pruebas de aceptacion}
\label{tab:aceptacion}
\begin{tabular}{p{0.9cm}p{8.2cm}p{1.6cm}p{2.3cm}}
\toprule
\textbf{ID} & \textbf{Criterio} & \textbf{Cumple} & \textbf{Observacion} \\
\midrule
1 & El sitio educativo se despliega correctamente en Vercel sin errores visibles. & Pendiente & Completar. \\
2 & El proyecto carga en menos de 3 segundos desde Vercel. & Pendiente & Medir. \\
3 & El endpoint devuelve datos correctos. & Pendiente & Postman. \\
4 & La navegacion o interaccion del aplicativo es fluida. & Pendiente & Validacion manual. \\
5 & No se presentan errores visibles en consola del navegador. & Pendiente & Captura. \\
6 & Las pruebas unitarias pasan correctamente en el pipeline automatico. & Pendiente & GitHub Actions. \\
7 & Las pruebas de integracion con Postman son exitosas. & Pendiente & Captura. \\
8 & Los resultados de JMeter estan dentro de los tiempos aceptables. & Pendiente & Reporte. \\
\bottomrule
\end{tabular}
\end{table}

\section{Analisis de Resultados}

Los resultados deben analizarse una vez se ejecute el proyecto en local, GitHub Actions, Postman, JMeter y Vercel. La interpretacion debe relacionar los hallazgos con la usabilidad, la capacidad de aprendizaje y la operabilidad.

Se espera que el proyecto demuestre que una interfaz organizada por temas, con audio, elementos visuales, laboratorio guiado y evaluaciones cortas, mejora la comprension del contenido y facilita la navegacion autonoma de los estudiantes.

\section{Mejoras Propuestas}

\begin{itemize}
  \item Agregar autenticacion para guardar avance por estudiante.
  \item Incorporar videos propios producidos por el colegio.
  \item Agregar accesibilidad con contraste configurable y lectura de texto ampliada.
  \item Registrar resultados de evaluaciones para analitica academica.
  \item Crear mas modulos de Ciencias Naturales para otros grados.
\end{itemize}

\section{Conclusiones}

El proyecto \textit{Mentes Creativas: Ciencias Naturales} integra desarrollo frontend, aseguramiento de calidad, pruebas automatizadas, integracion continua y despliegue en la nube. La aplicacion cumple un objetivo educativo concreto y permite evidenciar practicas de calidad asociadas al modelo ISO/IEC 25010.

La seleccion de usabilidad como caracteristica principal resulta adecuada porque los usuarios son estudiantes de primaria. Por esta razon, los criterios de calidad se enfocan en claridad, aprendizaje, navegacion, retroalimentacion y disponibilidad del sistema.

\printbibliography

\end{document}
